\documentclass[10pt]{article}
\usepackage[T1]{fontenc}
\usepackage[utf8]{inputenc}
\usepackage[french]{babel}
\usepackage{amsmath,amssymb,mathrsfs,mathpazo}
\usepackage{geometry}
\geometry{a4paper,margin=2cm}
\usepackage{graphicx,array,tabularx,tabularray,colortbl,multirow,multicol,booktabs,enumitem,pifont,ulem,fancybox,diagbox,cancel}
\usepackage{pstricks,pstricks-add,pst-plot,pst-tree,pst-node,pst-text,pst-func,pst-math,pst-bspline,pst-3dplot,pst-eucl,pst-solides3d}
\let\pscalculate\undefined
\usepackage{pst-bezier}
\usepackage{tikz,pgfplots,tkz-tab}
\usetikzlibrary{arrows,arrows.meta,patterns,calc,positioning,trees,decorations.markings,shapes.geometric,shapes.misc,shapes.arrows,intersections,decorations.pathmorphing,decorations.pathreplacing,backgrounds,shadows,babel}
\pgfplotsset{compat=1.18}
\usepgfplotslibrary{fillbetween}
\usepackage[np]{numprint}
\usepackage{listings,makecell,pas-tableur,twoopt}
\let\barre\undefined
\usepackage[pstricks]{tablvar}
\definecolor{darkgreen}{RGB}{0,100,0}
\providecommand{\red}{\color{red}}
\providecommand{\blue}{\color{blue}}
\providecommand{\green}{\color{green}}
\providecommand{\black}{\color{black}}
\providecommand{\white}{\color{white}}
\providecommand{\gray}{\color{gray}}
\providecommand{\cyan}{\color{cyan}}
\providecommand{\magenta}{\color{magenta}}
\providecommand{\yellow}{\color{yellow}}
\providecommand{\violet}{\color{violet}}
\providecommand{\orange}{\color{orange}}
\providecommand{\euro}{\texteuro}
\providecommand{\EUR}{\texteuro}
\providecommand{\ieme}{\textsuperscript{e}}
\providecommand{\degres}{\ensuremath{^\circ}}
\providecommand{\degree}{\ensuremath{^\circ}}
\providecommand{\vv}[1]{\overrightarrow{#1}}
\providecommand{\Coord}[1]{#1}
\providecommand{\decosix}{}
\providecommand{\starredbullet}{\ensuremath{\bullet}}
\pagestyle{empty}
\setlength{\parindent}{0pt}
\usepackage{adjustbox}\DeclareUnicodeCharacter{2236}{:}\DeclareUnicodeCharacter{2190}{\ensuremath{\leftarrow}}
\providecommand{\R}{\mathbb{R}}
\providecommand{\N}{\mathbb{N}}
\providecommand{\Z}{\mathbb{Z}}
\providecommand{\Q}{\mathbb{Q}}
\providecommand{\D}{\mathbb{D}}
\providecommand{\C}{\mathbb{C}}
\providecommand{\e}{{\mathrm{e}}}
\providecommand{\vect}[1]{\overrightarrow{#1}}
\providecommand{\vectt}[1]{\overrightarrow{\mathrm{#1}}}
\providecommand{\vvt}[1]{\overrightarrow{\mathrm{#1}}}
\providecommand{\np}[2][]{{#2~#1}}
\providecommand{\vv}[1]{{\overrightarrow{#1}}}
\renewcommand{\d}{{\mathrm{d}}}
\providecommand{\degres}{{^\circ}}
\providecommand{\degre}{{^\circ}}
\providecommand{\blue}{\color{blue}}
\providecommand{\red}{\color{red}}
\providecommand{\enskip}{\quad}
\providecommand{\llbracket}{\lbrack\!\lbrack}
\providecommand{\rrbracket}{\rbrack\!\rbrack}
\providecommand{\euro}{{\text{€}}}
\providecommand{\no}{{\text{n°}}}
\providecommand{\interval}[2]{{[#1;#2]}}
\providecommand{\intervFO}[2]{{[#1;#2[}}
\providecommand{\dotfill}{\ldots}

\providecommand{\ds}{\displaystyle}
\providecommand{\textding}[1]{#1}
\providecommand{\Oij}{\(\mathrm{O};\vec{\imath},\vec{\jmath}\)}
\providecommand{\Oijk}{\(\mathrm{O};\vec{\imath},\vec{\jmath},\vec{k}\)}
\providecommand{\Ouv}{\(\mathrm{O};\vec{u},\vec{v}\)}
\providecommand{\intervFF}[2]{\ensuremath{[#1;#2]}}
\renewcommand{\euro}{\ifmmode\text{\texteuro}\else\texteuro\fi}
\DeclareUnicodeCharacter{00B0}{\ifmmode^{\circ}\else\textdegree\fi}
\DeclareUnicodeCharacter{20AC}{\ifmmode\text{\texteuro}\else\texteuro\fi}
% Énoncés : Éducation nationale. Transcription et corrigés : APMEP, Denis Vergès et contributeurs. Découpage et classement : CoopMaths, Sébastien Lozano, Éric Elter et contributeurs. Adaptation de la présentation : Maths Explorer.
% CC BY-SA : https://coopmaths.fr/www/about/
% https://coopmaths.fr/alea/static/bac/2024/tex/bac_2024_09_sujet1_polynesie_1.tex
\begin{document}
\noindent\textit{APMEP / CoopMaths. CC BY-SA. Présentation adaptée par Maths Explorer.}\par\medskip
\medskip

Une concession automobile vend deux sortes de véhicules :


\begin{itemize}[label=$\bullet~$]
\item 60\,\% sont des véhicules tout-électrique ;
\item 40\,\% sont des véhicules hybrides rechargeables.
\end{itemize}

\medskip

75\,\% des acheteurs de véhicules tout-électrique et 52\,\% des acheteurs de véhicules hybrides ont la possibilité matérielle d'installer une borne de recharge à domicile.

On choisit un acheteur au hasard et on considère les évènements suivants:

\begin{itemize}[label=$\bullet~$]
\item $E$ : \og l'acheteur choisit un véhicule tout-électrique \fg{} ;
\item $B$ : \og l'acheteur a la possibilité d'installer une borne de recharge à son domicile \fg.
\end{itemize}

\medskip

\emph{Dans l'ensemble de l'exercice, les probabilités seront arrondies au millième si nécessaire.}

\medskip

\begin{enumerate}
\item Calculer la probabilité que l'acheteur choisisse un véhicule tout-électrique et qu'il ait la possibilité d'installer une borne de recharge à son domicile.

\emph{On pourra s'appuyer sur un arbre pondéré.}
\item Démontrer que $P(B) = 0,658$.
\item Un acheteur a la possibilité d'installer une borne de recharge à son domicile. Quelle
est la probabilité qu'il choisisse un véhicule tout-électrique ?
\item On choisit un échantillon de $20$ acheteurs. On assimile ce prélèvement à un tirage avec remise.

On note $X$ la variable aléatoire qui donne le nombre total d'acheteurs pouvant installer une borne de recharge à leur domicile parmi l'échantillon de $20$~acheteurs.
	\begin{enumerate}
		\item Déterminer la nature et les paramètres de la loi de probabilité suivie par $X$.
		\item Calculer $P(X = 8)$.
		\item Calculer la probabilité qu'au moins 10 acheteurs puissent installer une borne de recharge.
		\item Calculer l'espérance de $X$.
		\item La directrice de la concession décide d'offrir l'installation de la borne de recharge aux acheteurs ayant la possibilité d'en installer une à leur domicile. Cette installation coûte \np{1200}~\euro.
		
En moyenne, quelle somme doit-elle prévoir d'engager pour cette offre lors de la vente de $20$ ~véhicules?
	\end{enumerate}
\end{enumerate}

\end{document}
